\section{Related Work}

Decision-tree/Na\"ive-Bayes hybrids predate Farid et al.'s paper by nearly
two decades. Kohavi's NBTree~\cite{kohavi1996scaling} grows a tree and
places a local NB classifier at each leaf, using the tree to partition the
instance space rather than to select or weight attributes. Closer to
Algorithm 2's mechanism, Hall~\cite{hall2007decision} weights each
attribute for NB by a function of the depth at which a decision tree
tests it, on the same reasoning Farid et al.\ use (attributes near the
root separate classes more cleanly and should count more) but the two
papers were developed independently and Farid et al.\ do not cite
Hall's. Algorithm 1's filter-then-grow design continues a line that goes
back to John's robust decision
trees~\cite{john1995robust}, which also delete training instances a
classifier misclassifies before regrowing a tree on the remainder; the
family shares the same basic assumption that disagreement with a fitted
model is a reasonable proxy for label noise.

That assumption is the most contested part of the noise-filtering
literature. Brodley and Friedl~\cite{brodley1999identifying} show that a
single classifier used as its own noise judge is the weakest filtering
design available, because its errors on ambiguous but correctly labeled
instances are indistinguishable from its errors on genuinely mislabeled
ones; ensemble and consensus filters are more conservative and better
calibrated. Wong et al.~\cite{wong2020instance} build a more recent
instance-filtering hybrid that outperforms earlier one-shot filters
including designs close to Algorithm 1, and confident-learning
methods~\cite{northcutt2021confident} replace hard deletion with an
estimated confidence score per instance so uncertain cases can be
downweighted rather than discarded. On the correction side, Jiang et
al.~\cite{jiang2024correction} report that correcting a mislabeled
instance's label is often more effective than removing the instance
outright, which is a direct critique of Algorithm 1's delete-only design.
On the attribute side, tree-derived weighting for NB has continued past
Hall's original scheme; Zhang et al.~\cite{zhang2023multiview} build a
multi-view attribute-weighted NB from an ensemble of random trees rather
than a single tree, trading Algorithm 2's simplicity for a less
variance-prone weight estimate. Whether instance and attribute operations
should run in a fixed order, and whether that order matters, has direct
evidence in two places: Tsai et al.~\cite{tsai2013genetic} use a genetic
algorithm to learn a priority between feature and instance selection and
find the priority itself affects accuracy, and S\'aez et
al.~\cite{saez2015smoteipf} show that running a noise filter before
resampling strips borderline minority instances that resampling after
filtering would have kept, a sharp demonstration that pipeline order is
not a free choice. Both motivate directly testing the two possible orders
of Algorithms 1 and 2 rather than assuming they commute, which is what
the E9 study in Section~\ref{sec:experiments} does.

The third relevant body of work is methodological rather than about
classifiers. Cawley and Talbot~\cite{cawley2010overfitting} give the
general argument that any model-selection or preprocessing step fit on
data that later contributes to the reported test score produces an
optimistically biased estimate, and prescribe nesting that step inside
cross-validation; this is precisely the choice Farid et al.'s protocol
never states, and precisely what protocols A and B in this paper are
designed to separate. Demircio\u{g}lu~\cite{demircioglu2021measuring}
quantifies the resulting bias in a medical-imaging setting at up to 0.15
AUC or 0.17 accuracy when feature selection is fit outside
cross-validation, worse as the ratio of samples to features shrinks, which
matches the small-sample, many-attribute character of several datasets
used here. Kapoor and Narayanan~\cite{kapoor2023leakage} survey the same
failure mode across machine-learning-based science broadly and argue it
is common enough to be a contributor to the field's reproducibility
problems, and Bouke and Abdullah~\cite{bouke2023empirical} demonstrate the
effect concretely on NSL-KDD, one of the ten datasets used in this
replication, which is one reason NSL-KDD receives separate attention in
Section~\ref{sec:results}.
